% TFM-PROTECTED-TIKZ: fig:robot
% Figura fundamental: no eliminar, desactivar, rasterizar ni excluir del PDF.
\providecolor{VIUDeep}{HTML}{C63C32}

\begin{figure}[H]
\centering
\begin{tikzpicture}[>={Stealth[length=5pt]},font=\small,x=1cm,y=1cm]
  \draw[->,VIUGray!80] (0,0) -- (1.4,0) node[right] {$x$};
  \draw[->,VIUGray!80] (0,0) -- (0,1.3) node[above] {$y$};
  \node[VIUGray!80,below left=-1pt] at (0,0) {$O$};

  \def\cx{4.7}\def\cy{2.0}\def\ang{30}\def\rb{0.85}

  \draw[densely dotted,VIUGray] (\cx,\cy) -- (\cx+2.7,\cy);
  \draw[->,VIUGray] (\cx+1.55,\cy) arc (0:\ang:1.55);
  \node[VIUGray] at (\cx+1.82,\cy+0.42) {$\theta_i$};

  \begin{scope}[shift={(\cx,\cy)},rotate=\ang]
    \fill[VIUDark] (-0.14, \rb-0.02) rectangle (0.14, \rb+0.48);
    \fill[VIUDark] (-0.14,-\rb-0.48) rectangle (0.14,-\rb+0.02);
    \draw[VIUDark!75,thick] (0,-\rb) -- (0,\rb);
    \draw[very thick,draw=VIUOrange,fill=VIUOrange!12] (0,0) circle (\rb);
    \draw[->,thick,VIUOrange]
      ([shift=(55:0.5)]0,0) arc (55:160:0.5);
    \node[VIUOrange] at ([shift=(108:0.78)]0,0) {$\omega_i$};
    \draw[->,very thick,VIUDeep] (0,0) -- (\rb+1.0,0);
    \node[VIUDeep,anchor=south] at (\rb+0.55,0.05) {$v_i$};
    \draw[densely dashed,VIULinkBlue,thick]
      (\rb,0) -- (\rb+1.95,0);
    \fill[VIULinkBlue] (\rb+1.95,0) circle (2pt);
    \node[VIULinkBlue,anchor=south] at (\rb+1.95,0.07)
      {$\boldsymbol h_i$};
    \draw[<->,VIUGray] (0,-1.45) -- (\rb+1.95,-1.45);
    \node[VIUGray,fill=white,inner sep=1pt] at (1.4,-1.45) {$\ell$};
  \end{scope}

  \fill[VIUDark] (\cx,\cy) circle (1.8pt);
  \node[VIUDark,anchor=north east] at (\cx-0.05,\cy-0.08)
    {$(p_{x,i},p_{y,i})$};
\end{tikzpicture}
\caption[Robot móvil de tracción diferencial (uniciclo)]{Modelo geométrico del
robot diferencial. El estado es $q_i=(p_{x,i},p_{y,i},\theta_i)$ y las entradas
son la velocidad lineal $v_i$ y la angular $\omega_i$. El punto auxiliar
$\boldsymbol h_i=p_i+\ell(\cos\theta_i,\sin\theta_i)^\top$ explicita la
dirección de avance a una distancia $\ell$; no añade un estado físico.}
\label{fig:robot}
\viuownsource
\end{figure}
